\chapter{Conception et Architecture Système}

\section*{Introduction du chapitre}
Après avoir formalisé les besoins, ce chapitre traduit les exigences de QuickTriage en une
architecture logicielle cohérente. Nous présentons d'abord l'architecture générale en
couches, puis nous détaillons chacune des composantes : le pipeline de données et
d'apprentissage, le backend J2EE exposant le modèle, et le frontend mobile Flutter. Une
section est consacrée à la sécurité des données de santé, particulièrement sensible dans le
domaine médical. Le chapitre se conclut par les choix d'expérience et d'interface
utilisateur (UX/UI).

\section{Architecture générale du système}

\subsection{Vue d'ensemble en trois couches}
QuickTriage adopte une architecture en trois couches faiblement couplées, organisée autour
d'un principe directeur : \textbf{découpler l'entraînement du modèle de son déploiement}.
Cette séparation permet de faire évoluer indépendamment le modèle d'apprentissage, la logique
applicative et l'interface utilisateur. Les trois couches sont les suivantes :

\begin{enumerate}
  \item \textbf{La couche d'apprentissage automatique (\textit{offline}).} Réalisée en
  Python, elle prend en charge la préparation des données, l'entraînement, l'optimisation et
  la calibration du modèle. Son livrable est un artefact interopérable au format \textbf{ONNX}
  accompagné de ses métadonnées de prétraitement.
  \item \textbf{La couche backend (\textit{serveur d'application}).} Réalisée en Jakarta EE et
  déployée sur \textbf{WildFly~10}, elle charge le modèle ONNX, applique le même prétraitement
  qu'à l'entraînement, exécute l'inférence et expose le service via une API REST
  \textbf{JAX-RS}.
  \item \textbf{La couche frontend (\textit{client mobile}).} Réalisée en \textbf{Flutter},
  elle assure la saisie des constantes vitales, l'appel à l'API et la restitution visuelle du
  résultat au personnel soignant.
\end{enumerate}

\begin{infobox}
La frontière nette entre la phase \textit{offline} (entraînement Python) et la phase
\textit{online} (inférence Java) est rendue possible par le format ONNX, qui sérialise le
modèle de façon indépendante du langage. Le serveur Java n'a ainsi aucune dépendance vis-à-vis
de Python à l'exécution.
\end{infobox}

\subsection{Flux global de l'information}
Le cheminement d'une requête de triage illustre l'articulation des couches : le soignant
saisit les constantes dans l'application Flutter ; celles-ci sont sérialisées en JSON et
transmises par requête HTTP \texttt{POST} à l'API REST ; le backend désérialise la requête,
reconstruit le vecteur de caractéristiques, applique l'imputation et la standardisation, puis
exécute l'inférence ONNX ; le résultat (classe, confiance, probabilités, indicateur de
criticité) est renvoyé en JSON et affiché à l'écran avec un code couleur de priorité. Le
tableau~\ref{tab:couches} synthétise les responsabilités et technologies de chaque couche.

\begin{table}[h]
\centering
\renewcommand{\arraystretch}{1.4}
\begin{tabular}{p{3cm} p{5.5cm} p{5cm}}
\toprule
\textbf{Couche} & \textbf{Responsabilités} & \textbf{Technologies} \\
\midrule
Apprentissage (offline) & Préparation, entraînement, optimisation, calibration, export &
Python, pandas, scikit-learn, XGBoost, Optuna, SMOTE, ONNX \\
Backend (serveur) & Chargement du modèle, prétraitement, inférence, exposition REST &
Jakarta EE, WildFly 10, JAX-RS, EJB, ONNX Runtime \\
Frontend (mobile) & Saisie, validation, appel API, restitution visuelle &
Flutter, Dart, HTTP \\
\bottomrule
\end{tabular}
\caption{Responsabilités et technologies des trois couches de QuickTriage.}
\label{tab:couches}
\end{table}

\section{Architecture de données et pipeline d'apprentissage}

\subsection{Du jeu de données brut au modèle déployable}
La couche d'apprentissage suit un pipeline séquentiel dont chaque étape produit un artefact
exploitable par la suivante. Ce pipeline, conçu pour être reproductible, comprend :

\begin{enumerate}
  \item \textbf{Acquisition et diagnostic.} Chargement du jeu de données \textit{KTAS
  Emergency Triage} (1~267 patients réels) et contrôle de qualité, incluant la détection de
  toute variable susceptible d'introduire une fuite de données.
  \item \textbf{Nettoyage et prétraitement.} Traitement des valeurs manquantes par imputation
  (médianes), encodage des variables catégorielles et construction du vecteur de
  caractéristiques (23 \textit{features}).
  \item \textbf{Augmentation hybride.} Rééquilibrage des classes par combinaison de génération
  synthétique calibrée et de l'algorithme SMOTE, aboutissant à un jeu équilibré de 3~000
  patients (600 par classe).
  \item \textbf{Standardisation.} Normalisation des variables par \textit{StandardScaler}
  (centrage-réduction), dont les paramètres (moyennes, écarts-types) sont sauvegardés.
  \item \textbf{Entraînement et optimisation.} Apprentissage d'un modèle XGBoost dont les
  hyperparamètres sont optimisés par Optuna.
  \item \textbf{Calibration et évaluation.} Réglage du seuil de décision de la classe
  \textsc{Critique} et évaluation orientée sécurité patient (rappel des cas graves).
  \item \textbf{Export.} Sérialisation du modèle au format ONNX et externalisation des
  métadonnées de prétraitement dans un fichier JSON.
\end{enumerate}

\subsection{Externalisation des métadonnées de prétraitement}
Un choix de conception déterminant consiste à \textbf{externaliser les paramètres de
prétraitement} dans un fichier \texttt{metadata\_v3.json} plutôt que de les coder en dur. Ce
fichier contient les moyennes et écarts-types du \textit{StandardScaler}, les médianes de
l'imputeur, le seuil de la classe \textsc{Critique} (0,30) et la version du modèle. Le backend
le charge au démarrage, garantissant ainsi que \textbf{le prétraitement appliqué à l'inférence
est rigoureusement identique à celui de l'entraînement}. Cette cohérence est essentielle :
toute divergence de normalisation entre l'entraînement et l'inférence dégraderait silencieusement
les performances.

\begin{successbox}
L'externalisation des métadonnées dans \texttt{metadata\_v3.json} permet de mettre à jour le
modèle (nouvel ONNX + nouvelles métadonnées) \textbf{sans recompiler le backend}, améliorant
nettement la maintenabilité de la chaîne.
\end{successbox}

\section{Architecture du backend J2EE}

\subsection{Organisation en composants}
Le backend repose sur trois composants Jakarta EE aux responsabilités clairement séparées,
selon une logique de couches (présentation REST, logique métier, accès au modèle) :

\begin{itemize}
  \item \texttt{OnnxModelLoader} -- un bean \textbf{\texttt{@Singleton @Startup}} chargé une
  seule fois à l'amorçage du serveur. Il instancie l'environnement et la session ONNX Runtime
  et les maintient en mémoire pendant toute la durée de vie de l'application.
  \item \texttt{TriagePredictionService} -- un EJB \textbf{\texttt{@Stateless}} qui porte la
  logique métier : construction du vecteur de caractéristiques, imputation, standardisation,
  appel à la session ONNX et application du seuil de criticité.
  \item \texttt{TriageResource} -- une ressource \textbf{JAX-RS} (\texttt{@Path
  "/api/triage"}) qui expose le point d'entrée REST \texttt{POST /predict}, désérialise la
  requête JSON et sérialise la réponse.
\end{itemize}

\subsection{Justification des choix d'architecture}
Le recours au patron \textbf{Singleton} pour le chargement du modèle répond à une exigence de
performance : le chargement d'un modèle ONNX et l'initialisation du \textit{runtime} sont des
opérations coûteuses qu'il serait inacceptable de répéter à chaque requête. En les réalisant
une seule fois au démarrage, on garantit une inférence quasi instantanée par la suite. Le
caractère \textbf{\texttt{@Stateless}} du service métier favorise la scalabilité : le conteneur
peut instancier autant de beans que nécessaire pour traiter les requêtes concurrentes, sans
état partagé à synchroniser. Enfin, l'exposition via \textbf{JAX-RS} assure l'interopérabilité
et le découplage vis-à-vis du client mobile, conformément aux besoins non-fonctionnels.

\subsection{Contrat d'interface REST}
L'API expose un point d'entrée principal, décrit au tableau~\ref{tab:rest}. L'échange repose
sur le format JSON, standard et léger, adapté à une consommation mobile.

\begin{table}[h]
\centering
\renewcommand{\arraystretch}{1.4}
\begin{tabular}{p{2cm} p{4cm} p{7.5cm}}
\toprule
\textbf{Méthode} & \textbf{Chemin} & \textbf{Description} \\
\midrule
\texttt{POST} & \texttt{/api/triage/predict} & Reçoit les constantes vitales
(\texttt{TriageRequest}) et retourne la classe KTAS prédite, la confiance, les probabilités
par classe et l'indicateur de criticité (\texttt{TriageResponse}). \\
\bottomrule
\end{tabular}
\caption{Point d'entrée principal de l'API REST de QuickTriage.}
\label{tab:rest}
\end{table}

\section{Sécurité des données de santé}
Les données manipulées par QuickTriage relèvent de la catégorie des \textbf{données de santé},
particulièrement sensibles et soumises à des exigences réglementaires strictes (par exemple le
RGPD en Europe). La conception intègre donc plusieurs principes de protection, en cohérence avec
les recommandations \textbf{OWASP} :

\begin{itemize}
  \item \textbf{Minimisation des données.} Seules les constantes strictement nécessaires à la
  prédiction sont collectées ; aucune donnée d'identification directe du patient n'est requise
  pour l'inférence.
  \item \textbf{Chiffrement des échanges.} Les communications entre l'application mobile et le
  backend doivent emprunter le protocole \textbf{HTTPS}, afin de protéger les données en transit.
  \item \textbf{Contrôle d'accès.} L'accès à l'API doit être restreint au personnel autorisé,
  via un mécanisme d'authentification (jetons), conformément au principe du moindre privilège.
  \item \textbf{Validation des entrées.} Les données reçues sont systématiquement validées
  (bornes physiologiques, types) afin de prévenir les injections et les requêtes malformées.
  \item \textbf{Traçabilité.} L'historisation des décisions de triage (horodatage, version du
  modèle, décision finale) assure l'auditabilité du système.
\end{itemize}

\begin{warningbox}
La sécurité n'est pas une option dans un système médical. La protection des données en transit
(HTTPS), le contrôle d'accès et la validation des entrées constituent le socle minimal de toute
mise en production de QuickTriage.
\end{warningbox}

\section{Conception de l'expérience et de l'interface utilisateur}

\subsection{Principes directeurs}
L'interface de QuickTriage est conçue pour un usage \textbf{sous pression}, dans un
environnement où chaque seconde compte. Trois principes guident sa conception :
\begin{itemize}
  \item \textbf{Rapidité de saisie} : des champs ordonnés selon le flux naturel de recueil des
  constantes, des claviers numériques adaptés et des valeurs par défaut raisonnables.
  \item \textbf{Lisibilité immédiate du résultat} : un \textbf{code couleur de priorité}
  (rouge pour \textsc{Critique}, orange pour \textsc{Urgent}, vert pour \textsc{Non urgent})
  permettant une lecture instantanée, complété par le délai maximal et le niveau de confiance.
  \item \textbf{Confiance et transparence} : l'affichage des probabilités par classe et de la
  version du modèle donne au soignant les éléments pour exercer son jugement critique sur la
  proposition.
\end{itemize}

\subsection{Code couleur et hiérarchie visuelle}
Le code couleur reprend les conventions du triage hospitalier afin de ne pas dérouter le
personnel. La hiérarchie visuelle met en avant l'information la plus critique (la classe et,
le cas échéant, l'alerte de criticité), reléguant les détails (probabilités, métadonnées) à un
second plan consultable à la demande. Ce parti pris réduit la charge cognitive et accélère la
décision.

\section*{Conclusion du chapitre}
Ce chapitre a présenté l'architecture en trois couches de QuickTriage et justifié ses choix
structurants : découplage entraînement/déploiement via ONNX, externalisation des métadonnées
de prétraitement, composants J2EE aux responsabilités séparées, sécurisation des données de
santé et interface orientée terrain. Cette architecture constitue le squelette de la solution.
Le chapitre suivant en propose une vue formelle au travers de la modélisation UML : cas
d'utilisation, diagramme de classes et scénarios de séquence.
